\section{A bounded-ratio criterion for all sixteen capacities}

\begin{proposition}\label{prop:ratio}
If all sixteen capacities are positive and
\[
 \frac{\max_vw_v}{\min_vw_v}\leq r_*:=\frac{24+5\sqrt{21}}{17},
\]
the Clebsch blow-up has a fractional half of cost at most $N^2/50$,
and the integer conclusion follows for integer capacities.
\end{proposition}
\begin{proof}
Write $m=\min_vw_v$, $M=\max_vw_v$, and $r=M/m$.
Consider any matching side, with pair totals $s_1,\ldots,s_4$,
total mass $L$ and full cost $Q$. Then
\[
 Q\leq\frac14\sum_js_j^2,\qquad
 \frac{Q}{L^2}\leq\frac{\sum_jt_j^2}{4(\sum_jt_j)^2},
 \quad t_j=s_j/(2m)\in[1,r].
\]
Holding the other coordinates fixed, and writing their sum and sum of
squares as $S,A$, gives derivative
\[
 \frac{d}{dt}\frac{t^2+A}{(t+S)^2}
 =\frac{2(tS-A)}{(t+S)^3}.
\]
An interior turning point is a minimum. Each coordinate can therefore
be moved to an endpoint without decreasing the ratio, so its maximum
over the box occurs at a corner. With $k$ coordinates equal to $r$,
the desired bound $Q/L^2\leq2/25$ follows from
\[
 P_k(r)=8(kr+4-k)^2-25(kr^2+4-k)\geq0.
\]
Here $P_0=28$, $P_4=28r^2$, and
\[
 P_1=-17r^2+48r-3,\quad
 P_2=P_1+(r-1)(15-r),\quad
 P_3=P_1+14(r^2-1).
\]
For $1\leq r\leq r_*<3$, all are nonnegative: $r_*$ is the larger
root of $P_1$ and its smaller root is below $1$.
Choose the heavier side of any generator cut. Its mass $L\geq N/2$,
so proportional selection gives cost at most
$(N/(2L))^2(2L^2/25)=N^2/50$.
Apply Lemma~\ref{lem:rounding} for integer capacities.
\end{proof}

The constant $r_*\approx2.75958108675$ is sharp for the particular
four-pair product inequality: one pair with both capacities $M$ and
three pairs with both capacities $m$ attain its limiting equality.
This is not a sharpness claim for the sparse-half problem itself.
The existence of some neighbourhood of the uniform weighting already
follows from its strict slack and continuity; no novelty is asserted
for that elementary observation or for the present explicit criterion.
